\begin{theorem}[The arity-two equations on two legs and the uniform coefficients]
\label{thm:c11-oddar2}
With the solutions of Theorem \ref{thm:c9-quadratic} (iii), Theorem \ref{thm:c9-P1} (i) and Proposition \ref{prop:c11-oddsol}, write $\theta^\Delta_d$ for $\theta,\theta_\lambda,\theta_V$, $\theta^P_d$ for $\theta_P,\theta_{P,\lambda},\theta_{P,V}$, and $\theta^{\xi_\pm}_d$, $\theta^{\eta_\pm}_d$ for the coefficients of Proposition \ref{prop:c11-oddsol} along $\xi_\pm\in S_+\otimes W,\ S_-\otimes W^*$ and $\eta_+\in S_+\otimes W^*$, $\eta_-\in S_-\otimes W$, with $d\in\{V,\perp\}$ along the odd generators and $L=\log\tfrac32$. The classes of the arity-two residuals of Lemma \ref{lem:c9-sform} on two legs, the obstructions being the sums of the coefficients of $\varepsilon^0$ of the cycles with one special edge and two vertices $S^{(1)}(x)$, $S^{(1)}(y)$ or one vertex $S^{(2)}(x,y)$, vanish exactly under the conditions of Table \ref{tab:c11-oddar2}.

\begin{tabular}{l|l}
pair & condition \\
\hline
$(\Delta,P_b),\ (\Delta,\xi_\pm)$ & $\theta^\Delta_\perp=\theta^\Delta_\lambda=\theta^\Delta_V=-\tfrac13$ \\
$(\xi_+,\xi_-)$ & $\theta^P_d=\theta^{\xi_+}_V+\theta^{\xi_-}_V-\tfrac{27}4-2L\ \text{for every }d,\ \theta^{\xi_+}_V+\theta^{\xi_-}_V=\theta^{\xi_+}_\perp+\theta^{\xi_-}_\perp$ \\
$(\xi_\pm,\xi'_\pm)$ & $\theta^{\xi_\pm}_V=\theta^{\xi_\pm}_\perp$ \\
$(\xi_\pm,P_b)$ & $\theta^P_V=\theta^P_\perp=\theta^P_\lambda$ \\
$(\xi_+,\eta_-)$ & $\theta^{\xi_+}_V-\theta^{\xi_+}_\perp=\theta^{\eta_-}_V-\theta^{\eta_-}_\perp$ \\
$(\xi_+,\eta_+)$ & $\theta^R_\perp=\theta^R_\lambda=:\tau,\ \theta^\Delta_d=\theta^M_d=\tau-\tfrac13,\ \theta^{\xi_+}_d+\theta^{\eta_+}_d=\tau+\tfrac{77}{12}+2L\ (d=V,\perp)$ \\
\end{tabular}
\tabcaption{The conditions under which the classes of the two-leg arity-two residuals vanish, $L=\log\tfrac32$.\label{tab:c11-oddar2}}

On $(\xi_\pm,P_b)$ the obstruction is $\tfrac14\log\tfrac43$ times a $Q$-exact functional, and the arity-two generating cochain of two odd generators is $I+cS^{(1)}(\xi)+c'S^{(1)}(\xi')-cc'S^{(2)}(\xi,\xi')$ for odd parameters $c$, $c'$. Consequently:
\begin{itemize}
\item (i) every sector coefficient along $P_b$, $\Sigma$, $\Sigma^*$, $\mathfrak{so}(V)$ and $\mathfrak{so}(W_{\mathbb R}^\perp)$ is uniform, $\theta^z_d=:t_z$, and
\begin{equation}
t_\Delta=-\tfrac13,\qquad t_M=t_\Delta,\qquad t_R=t_\Delta+\tfrac13=0,\qquad t_{\xi_+}+t_{\xi_-}-t_P=\tfrac{27}4+2\log\tfrac32,\qquad t_{\xi_+}+t_{\eta_+}-t_\Delta=\tfrac{27}4+2\log\tfrac32;
\label{eq:c11-trel}
\end{equation}
\item (ii) the solutions form an affine space over the inner direction $t_z\mapsto t_z+c\lambda_z$, $\lambda_z=-1,+1,-2,+2,0$ on $\Sigma$, $\Sigma^*$, $P_b$, $K_b$ and the even rotations and $\Delta$, generated by the field redefinition $\phi^{(0)}=c\,\text{units}\cdot S^{(1)}(\Delta)$, and the outer direction $t_{\xi_\pm}\mapsto t_{\xi_\pm}\pm a$, $t_{\eta_\pm}\mapsto t_{\eta_\pm}\mp a$, the direction $G_B$ of Definition \ref{def:c8-B}; the relations \eqref{eq:c11-trel} and the sector uniformity are invariant under both, and they are the scheme-independent content of the computed pairs on the action side;
\item (iii) in the translation-covariant scheme of Definition \ref{def:c11-free}, $t_P=t_\Delta=-\tfrac13$, hence $t_{\xi_+}+t_{\xi_-}=\tfrac{77}{12}+2\log\tfrac32$ and $t_{\eta_+}=t_{\xi_-}$;
\item (iv) the individual $t_{\xi_\pm}$, $t_{\eta_\pm}$, $t_P$, the coefficients $\theta^K_d$, $\kappa^K_j$, $\kappa_V$, $\kappa_\lambda$ and every particular part with $\log\tfrac32$ are scheme artefacts.
\end{itemize}
The relations \eqref{eq:c11-trel} are the one-loop parts of $[\Delta,P_b]=-2P_b$, $[\Delta,\xi]=-\xi$, $[\xi,P_b]=0$, $[\xi_+,\xi_-]=P_{\Gamma(\xi_+,\xi_-)|_V}$ and of the bracket $[\xi,\eta]$ with its components along $\Delta$, $\mathfrak{so}(V)$ and $\mathfrak{so}(W_{\mathbb R}^\perp)$ (Proposition \ref{prop:c2-special}): with the corrected generators $S^{(1)}(z)+\hbar t_z\,\text{units}\cdot S^{(1)}(z)$ the arity-two equation along $(x,y)$ states that the order-$\hbar$ part of $(-1)^{|x|}\{S^{(1)}(x),S^{(1)}(y)\}+S^{(1)}([x,y])$ closes up to $\Lambda^{(0)}$-exact terms once the computed class of the obstruction is included.
\end{theorem}

\begin{proof}
In \eqref{eq:c9-arity2} the free directions enter through $F_1^{(1)}([x,y])$ and the brackets $\{F_1^{(1)}(x),S^{(1)}(y)\}$, $\{S^{(1)}(x),F_1^{(1)}(y)\}$; for uniform coefficients, $(-1)^{|x|}\{S^{(1)}(x),S^{(1)}(y)\}=-S^{(1)}([x,y])-\Lambda^{(0)}S^{(2)}(x,y)$ (Lemma \ref{lem:c2-components}) gives the change $(t_x+t_y-t_{[x,y]})$ units times $S^{(1)}([x,y])$ modulo exact terms, $t_{[x,y]}S^{(1)}([x,y])$ meaning the sum over the components of $[x,y]$ in the blocks. On two legs the sector parts enter separately through the brackets $\{S^V_2(\xi),S_2(\xi')\}$, $\{S_2(\xi),S^V_2(\xi')\}$ and their analogues, whose classes in the cohomology of $Q$ are combinations of the classes of the sector parts of $S^{(1)}([x,y])$ by the Fierz identities of the arrows $\lambda\to A|_V\to\lambda$ and $\lambda\to A|_{W_{\mathbb R}^\perp}\to\lambda$. The obstructions are the coefficients of $\varepsilon^0$ of the cycles of the statement, evaluated in two independent ways, and their classes are computed on complete bases of the two-leg functionals with the quantum numbers of the pair, the generating cochain of two odd generators having the sign forced by the classical master equation, $QS^{(2)}(\xi,\xi')+S^{(1)}(P_{\Gamma(\xi,\xi')|_V})-\{S^{(1)}(\xi),S^{(1)}(\xi')\}=0$ on two legs; the vanishing of the residual classes gives the table, the rows along $\Delta$ being Theorem \ref{thm:c9-theta2}. (i) $(\xi_\pm,\xi'_\pm)$ and $(\xi_\pm,P_b)$ give the uniformity along $\Sigma$ and $P_b$, then $(\xi_+,\eta_-)$ the uniformity along $\Sigma^*$, and $(\xi_+,\eta_+)$ with $\theta^\Delta_d=-\tfrac13$ gives $\tau=0$, $t_M=t_\Delta$, $t_R=0$ and $t_{\xi_+}+t_{\eta_+}=\tfrac{77}{12}+2L=t_\Delta+\tfrac{27}4+2L$; $(\xi_+,\xi_-)$ gives the fourth relation. (ii) The field redefinition generated by $\phi^{(0)}$ changes $F_1^{(1)}(x)$ by $\{\phi^{(0)},S^{(1)}(x)\}=c\lambda_x$ units times $S^{(1)}(x)$ modulo exact terms (Lemma \ref{lem:c2-components} with $S^{(2)}(\Delta,x)=0$), and $G_B$ is an $s$-closed cochain whose arity-one value along $\xi_\pm$ is $\pm S^{(1)}(\xi_\pm)$ and along $\eta_\pm$ is $\mp S^{(1)}(\eta_\pm)$ (Theorem \ref{thm:c8-obstruction} (iv)); both preserve the solvability of every equation, $\lambda$ and $B$ being derivations of $\mathfrak{psl}(4|4)$. The left sides of \eqref{eq:c11-trel} are invariant: the inner direction shifts $t_{\xi_\pm}$, $t_{\eta_\pm}$, $t_P$ by $-c$, $+c$, $-2c$ and leaves $t_\Delta$, $t_M$, $t_R$ fixed, and the outer direction shifts $t_{\xi_\pm}$ by $\pm a$ and $t_{\eta_\pm}$ by $\mp a$. (iii) The position part of $\hbar F_1^{(1)}(\Delta)$ along $b$ is $2\hbar F_1^{(1)}(P_b)$, so $t_P=t_\Delta$, and \eqref{eq:c11-trel}. (iv) Under the inner and outer directions the individual coefficients move.
\end{proof}
